\section{Implementazione}

Il progetto e' implementato in Python con una struttura modulare.
Le librerie principali sono NumPy, pandas, SciPy, scikit-learn, matplotlib e PyTorch.

\subsection{Struttura del codice}

\begin{lstlisting}[style=pythonstyle, caption={Struttura principale del progetto.}]
src/option_pricing_nn/
    black_scholes.py   # formula analitica e payoff
    monte_carlo.py     # simulazione Monte Carlo efficiente
    dataset.py         # generazione e salvataggio dati sintetici
    model.py           # rete feed-forward PyTorch
    training.py        # training loop, scaling, early stopping
    evaluation.py      # metriche di regressione
    plots.py           # visualizzazioni
    pipeline.py        # esperimento end-to-end
scripts/
    run_experiment.py  # entry point da linea di comando
\end{lstlisting}

\subsection{Efficienza computazionale}

La formula Black-Scholes e la generazione del dataset sono implementate in forma vettorializzata con NumPy.
Il metodo Monte Carlo usa la soluzione esatta del moto browniano geometrico al tempo $T$, evitando la discretizzazione temporale quando non necessaria.
Inoltre, la simulazione Monte Carlo e' divisa in batch sia sulle opzioni sia sui path simulati, in modo da controllare l'uso di memoria.

Il training della rete usa \texttt{DataLoader} di PyTorch, mini-batch gradient descent, standardizzazione degli input e scaling del target.
Queste scelte rendono il codice adatto sia a esperimenti rapidi sia a run piu' estesi.

\subsection{Riproducibilita'}

Per rendere gli esperimenti riproducibili, il progetto fissa i random seed di NumPy, Python e PyTorch.
Gli artefatti prodotti da un esperimento includono:

\begin{itemize}
    \item dataset sintetico generato;
    \item checkpoint del modello;
    \item scaler degli input;
    \item scaler del target;
    \item metriche in formato JSON;
    \item grafici diagnostici.
\end{itemize}

